\begin{abstract}
MiniMax Music 3 plans a track as eight residual vector quantised (RVQ) code streams at 25 Hz and renders them with a diffusion transformer over the latents of a continuous audio autoencoder. The encoder that maps audio back to those codes was never released, so a reference recording cannot condition the model. We learn a functional surrogate of that missing map by \emph{reverse distillation}: we generate 2,972 tracks with the model itself, keep the sampled codes together with the generator's top-50 code distributions, and train an encoder from the frozen autoencoder latents to the codes. Two design choices account for most of the result. First, the generator's code timeline is not a fixed ratio of the latent timeline but is stitched from 200-frame rollout chunks, and we recover the exact frame-to-latent correspondence and pool latents with it. Second, the acoustic codebooks are conditionally dependent across depth; replacing seven independent readouts by a decoder that is causal across codebook depth raises the downstream condition-replay cosine on 130 held-out generated tracks from 0.7698 to 0.8748 (paired bootstrap interval of the gain [0.101, 0.109]; community baseline 0.6633, true codes 0.9999) while exact-token agreement barely moves, suggesting substantial functional redundancy in the code space.

We release four encoders of 41M to 169M parameters, trained under maximal update parametrisation, together with the corpus. The recommended encoder reproduces its published token metrics to within 0.001 when re-evaluated on the same held-out split and encodes six minutes of stereo audio in 5.6 s on one RTX 3090. All evaluation is on model-generated audio and stops before the diffusion stage; calibration of the metric against random codes and a study on recorded music are the open validations.
\end{abstract}
